\documentclass[11pt]{article}\usepackage[margin=1in]{geometry}\usepackage{amsmath,amssymb,amsthm}\newtheorem{theorem}{Theorem}\newtheorem{remark}{Remark}
\title{Convex-Chart Geometry of Nonlinear Staple Dispersion}\author{}\date{21 July 2026}\begin{document}\maketitle
\begin{abstract}We compare the Fr\'echet dispersion of an $SU(2)$ staple family with a pairwise gauge-invariant proxy and quantify the complete curvature loss on the certified local chart.\end{abstract}
Let $S_1,\ldots,S_m$ lie in a geodesic ball of radius $r<\pi/2$ about their Fr\'echet mean $C$, and put
\[D=\sum_i d(C,S_i)^2,\qquad P={1\over m}\sum_{i<j}d(S_i,S_j)^2,\qquad s_r={\sin r\over r}.\]
\begin{theorem}If the Karcher equation $\sum_i\log_C S_i=0$ holds, then
\[s_r^2D\le P\le D.\]
Both quantities are invariant under a common endpoint gauge action.\end{theorem}
\begin{proof}Write $z_i=\log_C S_i$. Convexity of the ball and the differential singular values of the spherical exponential give
$s_r\|z_i-z_j\|\le d(S_i,S_j)\le\|z_i-z_j\|$. The Euclidean identity $m^{-1}\sum_{i<j}\|z_i-z_j\|^2=\sum_i\|z_i\|^2$ follows from $\sum_i z_i=0$. Summing proves the claim. Bi-invariance proves gauge invariance.\end{proof}
At the certified three-link staple radius $r=0.01043241535$, the factor $s_r^2$ exceeds $0.9999$, so convex-chart geometry causes negligible loss.
\begin{remark}This theorem is geometric. A normalized density comparison uniform over retained Yang--Mills boundary states is still required.\end{remark}\end{document}
